\centerline{\bf \huge 7 класс}
\centerline{\bf \huge Решения}

\problem{7.1}
В шестизначном числе пронумеровали слева направо все цифры числами от 1 до 6 , после чего все цифры с чётными номерами в том же порядке переставили на первые три места, а все цифры с нечётными номерами в том же порядке переставили на последние три места. Получившееся число совпало с исходным. Какое наибольшее количество различных цифр может быть в таком числе?

\textit{Ответ: 2.}

\textit{Решение.}
Обозначим это шестизначное число через $\overline{a_1 a_2 a_3 a_4 a_5 a_6}$. Тогда по условию оно совпадает с числом $\overline{a_2 a_4 a_6 a_1 a_3 a_5}$. Значит, $a_1=a_2=a_4$ и $a_3=a_6=a_5$. Таким образом, различных цифр не более двух. В качестве примера числа с двумя различными цифрами подойдёт число 112122.

\problem{7.2}
Найдите наименьшее число, начинающееся с цифр 1234567 и делящееся на 225.

\textit{Ответ: 1234567125.}

\textit{Решение.}
Искомое число должно делиться на 9 и на 25 , то есть оканчиваться на 00 или на 25 или на 50 или на 75 , и сумма цифр числа должна делиться на 9 . У числа 1234567 сумма цифр дает остаток 1 при делении на 9, то есть нам нужно приписать справа такие цифры, чтобы они давали число, делящееся на 25 и с суммой цифр, дающей остаток 8 при делении на 9. Наименьшее такое число - это 125 (дописать одну или две цифры при данных ограничениях невозможно).

\problem{7.3}
В банке можно положить вклад на месяц под фиксированный процент. В конце месяца банк начислит проценты, а потом округлит (по обычным правилам) начисленную сумму до целого числа рублей. У Васи и Пети одинаковые суммы денег. Вася разделил свои деньги на две равные части и открыл два вклада, а Петя положил все свои деньги на один вклад. Может ли оказаться, что после выплаты процентов сумма у Васи будет меньше суммы у Пети?

\textit{Ответ: может.}

\textit{Решение.} 
Пусть после начисления процентов у Васи на каждом из вкладов начислилось (с процентами) 1000,4 рублей, тогда после округления на каждом из двух его вкладов окажется 1000 рублей, то есть всего у Васи будет 2000 рублей. У Пети должна начислиться сумма $2 \cdot 1000,4=2000,8$ рублей. После округления у него будет 2001 рубль.

\problem{7.4}
За круглый стол сели семь математиков. Каждому из них дали карточку, на которой написано число «1» или « $-1 »$. Затем на доску записали семь чисел - произведения чисел на карточках каждых двух математиков, сидящих рядом. После этого на доску записали восьмое число - произведение всех семи чисел на карточках математиков. Какое наименьшее значение может принимать сумма 8 выписанных на доске чисел?

\textit{Ответ: -6.}

\textit{Решение.}
На доске будут записаны числа 1 или -1 . Нам нужно найти наибольшее возможное количество чисел -1 на доске. Заметим, что у каких-то двух сидящих рядом математиков на карточках будут одинаковые числа (иначе числа 1 и -1 чередовались бы по кругу, но это невозможно, так как 7 - нечётное число). Поэтому максимальное количество чисел -1 на доске будет 7 (а наименьшая возможная сумма - это -6 ). Покажем, что это возможно. Пусть у математиков (по кругу) будут карточки $1,-1,1,- 1,1,-1,1$. Так как произведение этих 7 чисел равно -1 , на доске будут выписаны 7 чисел -1 и одно число 1 .

\problem{7.5}
На острове живут рыцари и лжецы, рыцари всегда говорят правду, а лжецы всегда лгут. На чаепитие в зале собралось 30 человек. В процессе беседы десять из них сказали, что в зале ровно 20 рыцарей. Пятеро сказали, что в зале не больше 7 рыцарей. Двенадцать человек сказали, что в зале ровно 10 рыцарей. А остальные трое сказали, что в зале ровно 9 рыцарей. Каждый произнес ровно одну фразу. Сколько рыцарей присутствует на чаепитии?

\textit{Ответ: 5.}

\textit{Решение.}
Пусть в зале 10 рыцарей. Тогда те 12 человек, которые сказали, что в зале ровно 10 рыцарей, сказали правду, то есть они являются рыцарями. Получили противоречие, значит все эти 12 человек - лжецы. Но тогда рыцарей не больше чем $30- 12=18$ человек. Значит те, кто сказал, что в зале 20 рыцарей - тоже лжецы. Это сказали 10 человек, значит рыцарей не больше, чем $18-10=8$. Отсюда следует, что те 3 , кто сказал, что рыцарей 9 - тоже лжецы. Значит рыцарей не больше $8-3=5$. Значит, те 5 человек, кто сказал, что рыцарей не больше 7 , сказали правду. Поэтому они рыцари. А так как рыцарей не больше 5 , в комнате ровно 5 рыцарей.